%%%%%%%%%%%%%%%%%%%%%%%%%%%%%%%%%%%%%%%%%%%%%%%%%%%%%%%%%%%%%%%%%%%%%%%%%%%%%
% Schellekens 71 DW anomaly and HT descent criterion.
%
% Records the stratified inputs needed after the VOA classification:
% Type A (24 pure Niemeier, no finite orbifold), Type B (1 Leech Z/2),
% Type C (46 Leech Z/n, cyclic-orbifold presentations). Cyclic
% presentations are VOA input; HT descent also needs a trivialized
% Dijkgraaf-Witten class and finite-orbifold BV descent.
%%%%%%%%%%%%%%%%%%%%%%%%%%%%%%%%%%%%%%%%%%%%%%%%%%%%%%%%%%%%%%%%%%%%%%%%%%%%%

\chapter[Schellekens 71 $\alpha$ criterion]{The Schellekens 71
$\alpha$-criterion: finite-orbifold descent across the holomorphic
$c = 24$ landscape}
\label{ch:schellekens-71-alpha-classification-platonic}

\paragraph{Orientation.}
Chapter~\ref*{ch:monster-chain-level-e3-top-platonic} isolates the
Moonshine module $V^\natural$ as a finite-orbifold boundary
construction. The Leech $\mathbb{Z}/2$ presentation supplies a VOA;
the descended HT model requires the abelian Leech HT theory,
finite-orbifold BV descent, and the local datum
$[\alpha_{\mathrm{orb}}]=0$ in
$H^3(B\mathbb{Z}/2;\mathrm{U}(1))$. The equality
$\Lambda^\sigma=0$ follows from torsion-freeness and $\sigma=-1$;
Leech rootlessness identifies the weight-one-zero moonshine entry. It
does not by itself trivialize a finite-orbifold BV associator.

The cyclic-orbifold construction of van Ekeren--M\"oller--Scheithauer,
the Schellekens-list theorem of van Ekeren--Lam--M\"oller--Shimakura,
and the generalised deep-hole classification of
M\"oller--Scheithauer supply the $71$ VOA presentations. They do not
automatically produce original-complex $E_3$ or HT descent. This
chapter records the stratum-by-stratum criterion: Type~A is the
abelian lattice case; Types~B and~C require a trivialized
Dijkgraaf-Witten class and finite-orbifold BV descent for the chosen
cyclic action. The $\mathbb{Z}/3$ Coxeter--Todd case is a representative
Type~C local datum, not a replacement for the general BV hypothesis.

%%%%%%%%%%%%%%%%%%%%%%%%%%%%%%%%%%%%%%%%%%%%%%%%%%%%%%%%%%%%%%%%%%%%%%%%%%%%%
\section{Orbifold-type stratification recalled}
\label{sec:schellekens-71-stratification-recalled}
%%%%%%%%%%%%%%%%%%%%%%%%%%%%%%%%%%%%%%%%%%%%%%%%%%%%%%%%%%%%%%%%%%%%%%%%%%%%%

Recall that Schellekens 1993 gives the possible weight-one
Lie-algebra components $V_1$ of strongly rational holomorphic
VOAs of central charge $c = 24$. Seventy-one rows appear: a
list of reductive Lie algebras together with the single entry
$V_1 = 0$, for which we use the Moonshine representative
$V^\natural$. The uniform construction
theorem of van Ekeren--Lam--M\"oller--Shimakura (2020, Adv.~Math.;
\emph{Schellekens' list and the very strange formula}, arXiv:2005.12248)
establishes that every Schellekens VOA is realisable as a
$\mathbb{Z}/n$-orbifold of the Leech lattice VOA $V_\Lambda$ for an
automorphism $\sigma \in \mathrm{Co}_0 = \mathrm{Aut}(\Lambda)$ of some
finite order $n$. M\"oller-Scheithauer (2023, Ann.~Math.) refine this
into a bijection between generalised deep holes of $\Lambda$ and the
$70$ Schellekens entries with $V_1 \ne 0$.

\begin{definition}[Orbifold-type stratification]
\label{def:schellekens-orbifold-type}
For a Schellekens VOA $W$ with a Niemeier or Leech cyclic-orbifold
presentation from the cited classification theorems, with lattice
$\Lambda$ and automorphism
$g \in \mathrm{Aut}(V_\Lambda)$ of order $n$, the \emph{orbifold type}
of $W$ is:
\begin{itemize}
\item \textbf{Type A}: if $g = \mathrm{id}$ (pure Niemeier VOA);
\item \textbf{Type B}: if $\Lambda = \Lambda_{\mathrm{Leech}}$,
$n = 2$, and $\Lambda^\sigma = 0$;
\item \textbf{Type C}: if $\Lambda = \Lambda_{\mathrm{Leech}}$,
$n \ge 2$, and $\Lambda^\sigma \ne 0$.
\end{itemize}
For Types~B and~C the finite-orbifold BV obstruction is a
Dijkgraaf-Witten class
$[\alpha_{\mathrm{orb}}]\in H^3(B\mathbb{Z}/n;\mathrm{U}(1))$. Type~A
has no finite group to gauge and hence no such obstruction.
\end{definition}

\begin{remark}[Stratum counts, following M\"oller-Scheithauer 2023]
\label{rem:schellekens-strata-counts}
The stratification \emph{partitions} the $71$ Schellekens entries as
$24_{\mathrm{Niemeier}} + 1_{V_1=0} + 46_{\mathrm{nonlat}} = 71$:
\begin{itemize}
\item Type A: $24$ rows, one for each Niemeier lattice.
\item Type B: $1$ row, represented by
$V^\natural = V_{\Lambda_{\mathrm{Leech}}}^{\mathbb{Z}/2}$ with
$\sigma = -1$.
\item Type C: $46$ rows, corresponding to the $46$ generalised deep
holes of the Leech lattice of automorphism order $n \ge 2$ with
non-trivial fixed sublattice. The orders are distributed (by
M\"oller-Scheithauer 2023) as $n = 2$ ($15$ cases with
$\Lambda^\sigma \ne 0$), $n = 3$ ($7$ cases), $n = 4$ (several cases),
$n = 5, 6, 7, 8, 10, 11, 13, \ldots$ with individual entries, up to
maximum order $n \approx 46$.
\end{itemize}
There is no separate ``affine Kac-Moody combination'' tier: the
cyclic-orbifold construction theorem subsumes every such presentation into
Type~A (when the affine presentation agrees with a Niemeier lattice
VOA) or Type~C (when it is a true orbifold).
\end{remark}

%%%%%%%%%%%%%%%%%%%%%%%%%%%%%%%%%%%%%%%%%%%%%%%%%%%%%%%%%%%%%%%%%%%%%%%%%%%%%
\section{The descent criterion}
\label{sec:schellekens-71-main}
%%%%%%%%%%%%%%%%%%%%%%%%%%%%%%%%%%%%%%%%%%%%%%%%%%%%%%%%%%%%%%%%%%%%%%%%%%%%%

\begin{theorem}[Schellekens 71 finite-orbifold descent criterion;
\ClaimStatusConditional]
\label{thm:schellekens-71-all-alpha-zero}
Let $W$ be a chosen representative of one of the $71$ Schellekens rows
at central charge $c = 24$, with orbifold type as in
Definition~\ref{def:schellekens-orbifold-type}. Assume:
\begin{enumerate}[label=\textup{(\roman*)}]
\item \textbf{VOA presentation.} $W$ is identified by the Niemeier,
FLM, or cyclic-orbifold classification input in its stratum.
\item \textbf{BV anomaly datum.} The finite-orbifold BV obstruction is
trivialized in the non-vacuous strata:
 \begin{itemize}
 \item Type A ($24$): no finite orbifold is gauged.
 \item Type B ($1$): the Leech $\Z/2$ datum satisfies
 $\alpha_{\mathrm{orb}}(\sigma,\sigma,\sigma)=+1$.
 \item Type C ($46$): the chosen cyclic action is equipped with
 $[\alpha_{\mathrm{orb}}]=0$ in
 $H^3(B\mathbb{Z}/n;\mathrm{U}(1))$.
 \end{itemize}
\item \textbf{Finite-orbifold BV descent.} The corresponding
finite-orbifold BV construction applies to the abelian lattice HT
model and its boundary restriction identifies the descended boundary
chiral algebra with $W$.
\end{enumerate}
Then the descended HT factorisation algebra carries a chain-level
$E_3$-topological structure with boundary chiral algebra $W$. The
Schellekens classification supplies the VOA presentations; it does not
by itself supply the BV anomaly trivializations or the HT descent.
\end{theorem}

The verification is stratum-by-stratum.
Section~\ref{sec:schellekens-typeA-pure-niemeier} gives the vacuous
Type~A anomaly and abelian HT model,
Section~\ref{sec:schellekens-typeB-leech-z2} records the Type~B local
datum from Chapter~\ref*{ch:monster-chain-level-e3-top-platonic}, and
Section~\ref{sec:schellekens-typeC-level-matching} records the Type~C
cyclic-orbifold input and the remaining BV datum.

\begin{remark}[Stratified scope]
\label{rem:schellekens-stratified-scope}
The theorem is a stratified $c=24$ criterion. The partition
$24_{\mathrm{Niemeier}} + 1_{V_1=0} + 46_{\mathrm{nonlat}} = 71$ is
classification input; chain-level HT descent is obtained only after
the anomaly and BV-descent data in the theorem are supplied. The
Monster shortcut $\Lambda^\sigma = 0$ is available only in Type~B.
\end{remark}

%%%%%%%%%%%%%%%%%%%%%%%%%%%%%%%%%%%%%%%%%%%%%%%%%%%%%%%%%%%%%%%%%%%%%%%%%%%%%
\section{Type A: pure Niemeier lattice VOAs ($24$ cases)}
\label{sec:schellekens-typeA-pure-niemeier}
%%%%%%%%%%%%%%%%%%%%%%%%%%%%%%%%%%%%%%%%%%%%%%%%%%%%%%%%%%%%%%%%%%%%%%%%%%%%%

\begin{proposition}[Type A: vacuous DW class; \ClaimStatusProvedHere]
\label{prop:schellekens-typeA-alpha-zero}
Let $\Lambda$ be any of the $24$ Niemeier lattices (including the
Leech lattice). The pure lattice VOA $V_\Lambda$ is a strongly rational
holomorphic VOA of central charge $c = 24$ (FLM 1988, Ch.~2) and
appears on the Schellekens list as Type~A. Its orbifold presentation
is trivial ($g = \mathrm{id}$, $n = 1$), and
$\alpha_{\mathrm{orb}}(V_\Lambda) = 0$ in
$H^3(B\{1\}; \mathrm{U}(1)) = 0$.
\end{proposition}

\begin{proof}
For $g = \mathrm{id}$ there is no finite-group orbifold to gauge. The
Dijkgraaf-Witten class is therefore a class on the trivial group:
$H^3(B\{1\};\mathrm{U}(1))=0$. The determinant/sign formula for a
non-trivial cyclic orbifold is not the relevant computation in this
case.

Alternatively: a pure Niemeier VOA $V_\Lambda$ is the chiral algebra
of an abelian holomorphic Chern-Simons theory on $X \times \mathbb{R}$
with charge lattice $\Lambda$ (Costello-Li). This gauge theory is
anomaly-free at the level of 2-form symmetry because $\Lambda$ is
even unimodular: the pairing $\langle \cdot, \cdot \rangle \colon
\Lambda \otimes \Lambda \to \mathbb{Z}$ is unimodular, and
$\frac{1}{2} \langle \lambda, \lambda \rangle \in \mathbb{Z}$ gives
a well-defined level. No orbifold, no DW cocycle. Chain-level
$E_3$-topological structure follows from
Theorem~\ref*{thm:E3-topological-km} applied to the abelian case.
\end{proof}

\begin{remark}[Enumeration of the 24 Niemeier lattices]
\label{rem:schellekens-24-niemeier-enumeration}
The $24$ Niemeier lattices have root systems (by Niemeier 1973):
\begin{itemize}
\item $\Lambda_{\mathrm{Leech}}$ (root-free),
\item $A_1^{24}$, $A_2^{12}$, $A_3^8$, $A_4^6$,
$A_5^4D_4$, $A_6^4$, $A_7^2D_5^2$, $A_8^3$,
\item $A_9^2D_6$, $A_{11}D_7E_6$, $A_{12}^2$,
$A_{15}D_9$, $A_{17}E_7$, $A_{24}$,
\item $D_4^6$, $D_6^4$, $D_8^3$, $D_{10}E_7^2$,
$D_{12}^2$, $D_{16}E_8$, $D_{24}$,
\item $E_6^4$, $E_8^3$.
\end{itemize}
For a lattice VOA, $V_1$ contains the Heisenberg summand
$\Lambda \otimes_\Z \C$ of dimension~$24$, together with root vectors
from norm-$2$ lattice elements. Thus the Leech lattice VOA has
$V_1 \cong \C^{24}$ abelian; $V_1=0$ belongs instead to the Moonshine
orbifold $V^\natural$. All pure Niemeier lattice VOAs are Type~A: no
orbifold, vacuous DW class,
chain-level $E_3$-topological via abelian holomorphic CS on the charge
lattice (each Niemeier lattice gives an even unimodular
level-specification of abelian CS).
\end{remark}

%%%%%%%%%%%%%%%%%%%%%%%%%%%%%%%%%%%%%%%%%%%%%%%%%%%%%%%%%%%%%%%%%%%%%%%%%%%%%
\section{Type B: the Leech $\mathbb{Z}/2$ shortcut ($1$ case)}
\label{sec:schellekens-typeB-leech-z2}
%%%%%%%%%%%%%%%%%%%%%%%%%%%%%%%%%%%%%%%%%%%%%%%%%%%%%%%%%%%%%%%%%%%%%%%%%%%%%

\begin{proposition}[Type B: Leech $\Z/2$ anomaly datum;
\ClaimStatusConditional]
\label{prop:schellekens-typeB-alpha-zero}
The Moonshine module $V^\natural = V_{\Lambda_{\mathrm{Leech}}}^{\mathbb{Z}/2}$
at $\sigma = -1$ is the unique Type~B Schellekens entry. Its HT
descent requires the local finite-orbifold BV datum
\[
 \alpha_{\mathrm{orb}}(\sigma,\sigma,\sigma)=+1
\]
in $H^3(B\mathbb{Z}/2;\mathrm{U}(1))\cong\mathbb{Z}/2$, together with
finite-orbifold BV descent for the Leech HT model. The determinant
factor
$\mathrm{sign}(\det(1 + 1|_\Lambda)) = \mathrm{sign}(2^{24}) = +1$
is a compatibility check, not the whole anomaly computation.
\end{proposition}

\begin{proof}
See Chapter~\ref*{ch:monster-chain-level-e3-top-platonic},
Proposition~\ref*{prop:monster-alpha-explicit-zero} and
Theorem~\ref*{thm:monster-chain-level-e3-top}.
The Leech $\mathbb{Z}/2$ mechanism uses the simultaneous facts
(i) $\sigma = -1$ is an even-unimodular automorphism ($\sigma$-invariance
of the bilinear form is automatic because $\langle -\lambda, -\mu
\rangle = \langle \lambda, \mu \rangle$), and
(ii) the Leech lattice is torsion-free, so
$\Lambda^\sigma=\{\lambda: \lambda=-\lambda\}=0$. Rootlessness is used
to identify the resulting orbifold as the weight-one-zero moonshine
entry, not to prove a BV anomaly trivialization.
\end{proof}

\begin{remark}[What Type B does not generalise]
\label{rem:schellekens-typeB-non-generalisable}
The identity $\Lambda^\sigma=0$ for $\sigma=-1$ is a torsion-free
lattice fact; by itself it is not unique to the Leech lattice. What is
special is the Schellekens presentation: Type~B is the Leech
$\mathbb Z/2$ orbifold giving the unique weight-one-zero entry
$V^\natural$. Type~A has no orbifold anomaly to compute, and Type~C
uses cyclic automorphisms whose fixed lattices are generally
non-trivial. Thus the Type~B argument cannot be promoted to a uniform
proof of the $71$ entries.
\end{remark}

%%%%%%%%%%%%%%%%%%%%%%%%%%%%%%%%%%%%%%%%%%%%%%%%%%%%%%%%%%%%%%%%%%%%%%%%%%%%%
\section{Type C: Leech $\mathbb{Z}/n$ orbifolds with
$\Lambda^\sigma \ne 0$ ($46$ cases)}
\label{sec:schellekens-typeC-level-matching}
%%%%%%%%%%%%%%%%%%%%%%%%%%%%%%%%%%%%%%%%%%%%%%%%%%%%%%%%%%%%%%%%%%%%%%%%%%%%%

Type~C comprises the $46$ Schellekens entries of the form
$W = V_{\Lambda_{\mathrm{Leech}}}^g$ for $g \in \mathrm{Co}_0$ of
order $n \ge 2$ with $\Lambda^\sigma \ne 0$. These correspond
bijectively (M\"oller-Scheithauer 2023) to the $46$ generalised deep
holes of the Leech lattice with automorphism group of order $\ge 2$
and non-trivial fixed sublattice.

\subsection{Cyclic-orbifold level matching}

\begin{proposition}[Cyclic-orbifold presentation and level matching;
\ClaimStatusProvedElsewhere]
\label{prop:ems-level-matching}
Let $g \in \mathrm{Co}_0$ be of order $n$, with fixed sublattice
$\Lambda^\sigma \subseteq \Lambda_{\mathrm{Leech}}$ of rank $r$. The
cyclic-orbifold construction produces a strongly rational holomorphic
VOA at $c = 24$ only under the level-matching condition that the
conformal weight of the vacuum of the $g$-twisted sector
$V_{\Lambda_{\mathrm{Leech}}}^{\mathrm{tw}(g)}$ lies in
$\frac{1}{n}\mathbb{Z}$. This is the VOA orbifold input. The HT
descent theorem still requires a chosen trivialization of the
finite-orbifold BV class
$[\alpha_{\mathrm{orb}}]\in H^3(B\mathbb{Z}/n; \mathrm{U}(1))$.
\end{proposition}

\begin{proof}[Proof; cited]
The cyclic-orbifold existence and level-matching criterion are the
input from van Ekeren--M\"oller--Scheithauer,
\emph{Construction and classification of holomorphic vertex operator
algebras}. The Schellekens-list application and explicit
weight-one constraints are supplied by van Ekeren--Lam--M\"oller--
Shimakura, \emph{Schellekens' List and the Very Strange Formula},
together with M\"oller--Scheithauer,
\emph{Dimension formulae and generalised deep holes}, for the
generalised-deep-hole enumeration. The forward implication is the
modular-invariance requirement on the twisted character
$\mathrm{tr}_{V^{\mathrm{tw}(g)}}(g^i\, q^{L_0 - c/24})$; the
backward implication uses Dong-Li-Mason 2000
(\emph{Modular-invariance of trace functions in orbifold theory})
combined with Dong-Lepowsky-Mason 1996 simple-current theory. The
Dijkgraaf-Witten class is the finite-gauging obstruction of
Dijkgraaf-Witten 1990 and Dijkgraaf-Pasquier-Roche 1991; the HT
application uses its trivialization as a separate BV input.
\end{proof}

\begin{corollary}[Type C descent under local anomaly trivialization;
\ClaimStatusConditional]
\label{cor:schellekens-typeC-alpha-zero}
Every Type~C Schellekens VOA whose chosen cyclic action is equipped
with $[\alpha_{\mathrm{orb}}]=0$ in
$H^3(B\mathbb{Z}/n; \mathrm{U}(1))$ and finite-orbifold BV descent
admits the descended chain-level HT $E_3$ model of
Theorem~\ref{thm:schellekens-71-all-alpha-zero}.
\end{corollary}

\begin{proof}
Each Type~C entry $W = V_{\Lambda_{\mathrm{Leech}}}^g$ is, by
construction, a strongly rational holomorphic VOA of central charge
$c = 24$ with a cyclic-orbifold presentation satisfying the
level-matching input of Proposition~\ref{prop:ems-level-matching}. The
additional hypotheses are exactly the obstruction and descent
hypotheses in Theorem~\ref{thm:schellekens-71-all-alpha-zero}.
\end{proof}

\subsection{The $\mathbb{Z}/3$ Coxeter--Todd orbifold}
\label{subsec:schellekens-z3-coxeter-todd}

Take
$g \in \mathrm{Co}_0$ of order $3$, in the Frame-shape class
``$1^{-3} 3^9$'' (conjugacy class $3\mathrm{A}$ of $\mathrm{Co}_0$, the
unique class of order $3$ with $\Lambda^\sigma$ of rank $12$).

\begin{lemma}[Coxeter-Todd fixed sublattice]
\label{lem:coxeter-todd-leech-z3-fixed}
For $g \in \mathrm{Co}_0$ in the $1^{-3} 3^9$ class of order~$3$,
the fixed sublattice $\Lambda^\sigma$ is isomorphic to the Coxeter-Todd
lattice $K_{12}$: the unique even positive-definite lattice of rank
$12$, determinant $729=3^6$, and minimum norm $4$.
\end{lemma}

\begin{proof}[Proof; cited]
The isometry class of $\Lambda^\sigma$ for each $g \in \mathrm{Co}_0$
is determined by the cycle-shape classification (Conway 1969,
Curtis-Wilson tables). For Frame shape $1^{-3} 3^9$, the $\sigma$-fixed
sublattice has rank $12$. Matching lattice invariants (rank~$12$,
determinant~$3^6 = 729$, and minimum norm~$4$) identifies it as the
Coxeter-Todd lattice $K_{12}$ (Conway-Sloane 1988, \S 4.9). $\qed$
\end{proof}

\begin{proposition}[Coxeter--Todd $\mathbb{Z}/3$ anomaly datum;
\ClaimStatusConditional]
\label{prop:z3-coxeter-todd-alpha-zero-explicit}
For $g \in \mathrm{Co}_0$ of order~$3$ in class $3\mathrm{A}$
(Coxeter-Todd fixed sublattice), the cyclic-orbifold table gives the
local level-matching datum
$h_{\mathrm{tw}}=2\in\frac{1}{3}\mathbb{Z}$. If the local
Dijkgraaf-Witten/BV comparison identifies this datum with the
finite-orbifold associator, then $[\alpha_{\mathrm{orb}}]=0$ in
$H^3(B\mathbb{Z}/3; \mathrm{U}(1)) = \mathbb{Z}/3$. The compatibility
checks are: (i) the sign factor
$\mathrm{sign}(\det(1 - g|_\Lambda)) = +1$ because
the non-trivial eigenvalues come in complex-conjugate pairs
$(\zeta_3, \bar{\zeta_3})$ contributing
$|1 - \zeta_3|^2 = 3$ per pair, giving
$\det(1 - g|_{\Lambda^{\sigma,\perp}})=3^6>0$; (ii) the fixed lattice
is $K_{12}$, even of determinant $729=3^6$.
\end{proposition}

\begin{proof}
For the sign factor: the characteristic polynomial of
$g \in \mathrm{Co}_0$ on $\Lambda \otimes \mathbb{R}$ is
$t^{24} - 3 t^{21} + \cdots$ (Frame shape $1^{-3} 3^9$), which
factors as $(t - 1)^{-3}(t^3 - 1)^9$ formally, meaning $9$ copies of
the primitive cube roots of unity pair $(\zeta, \bar\zeta)$ with
$\zeta = e^{2\pi i/3}$ and $12$ fixed eigenvalues giving $\Lambda^\sigma$
(the correction from the $-3$ exponent accounts for the mismatch
between the cycle-shape and a real representation). On the orthogonal
complement $\Lambda^{\sigma,\perp}$ of rank~$12$, the operator $1 - g$
has eigenvalues $(1 - \zeta, 1 - \bar\zeta)^{\oplus 6}$, each pair
contributing $(1-\zeta)(1-\bar\zeta) = 1 - \zeta - \bar\zeta + 1
= 2 - (-1) = 3$ in absolute value of the product. Hence
$\det(1 - g|_{\Lambda^{\sigma,\perp}}) = 3^6 = 729 > 0$; sign factor
$+1$.

For the cocycle restriction: the FLM cocycle $\epsilon \colon \Lambda
\times \Lambda \to \{\pm 1\}$ is $\mathbb{Z}/2$-valued and encodes
the asymmetry of the VOA operator product at the underlying lattice
level. Restricting to $\Lambda^\sigma = K_{12}$, which is a sublattice
with $\mathbb{Z}/3$ action compatible with the inclusion, the
induced $3$-cocycle $\epsilon|_{K_{12}}$ in
$H^3(B\mathbb{Z}/3; \mathrm{U}(1))$ is the Kapustin-Saulina
anomaly representative.

The Dong-Mason quantum Galois theorem (J.~Alg.~214, 1999) plus the
DLM 2000 modular-invariance theorem supply the VOA modular constraint
on the twisted vacuum conformal weight. For the specific
$\sigma \in 3\mathrm{A}$ class, the cited cyclic-orbifold table records
$h_{\mathrm{tw}}(V_{\Lambda_{\mathrm{Leech}}}^{\mathrm{tw}(g)}) = 2$.
The integer~$2$ lies in $\frac{1}{3}\mathbb{Z}$, so level matching
holds. The equality $[\alpha_{\mathrm{orb}}]=0$ is the additional
local BV comparison datum stated in the proposition.

Chain-level: the orbifold Jacobi identity for $\mathbb{Z}/3$ is
established in Dong-Lepowsky-Mason 1996 (\emph{Simple currents}) +
Dong-Li-Mason 2000; combined with $\sigma$-equivariance of the lift
$\widetilde{\sigma}$ (fixed by the Dong-Mason uniqueness up to
$\mathrm{Hom}(\Lambda^\sigma / (1 - \sigma)\Lambda, \mathrm{U}(1))$,
and the level-matching condition fixes the VOA extension), the
orbifold product is a VOA chain identity. Chain-level
$E_3$-topological structure descends only after the finite-orbifold BV
descent and the local Dijkgraaf-Witten trivialization are supplied.
\end{proof}

\begin{remark}[Type C mechanism]
\label{rem:schellekens-typeC-universal-mechanism}
The Coxeter-Todd example of
Proposition~\ref{prop:z3-coxeter-todd-alpha-zero-explicit} is
representative of the local data one must supply: for every
$g \in \mathrm{Co}_0$ corresponding to a
generalised deep hole of the Leech lattice, the fixed sublattice
$\Lambda^\sigma$ is a well-identified even lattice (M\"oller-Scheithauer
2023 Table 3 enumerates all $46$), the sign factor
$\mathrm{sign}(\det(1 - g|_{\Lambda^{\sigma,\perp}}))$ is positive
(because non-trivial eigenvalues of $g$ on $\Lambda \otimes \mathbb{R}$
come in complex-conjugate pairs contributing positive absolute
values), and the $H^3$-cocycle restriction must be trivialized in the
finite-orbifold BV model. The Type~C entries differ in which lattice
$\Lambda^\sigma$ appears and what twisted-vacuum conformal weight
$h_{\mathrm{tw}}$ materialises. The listed entries are precisely the
level-matched cyclic orbifolds; the HT descent is the additional
finite-orbifold BV step.
\end{remark}

\begin{remark}[Explicit enumeration of Type C cases]
\label{rem:schellekens-typeC-enumeration}
Distribution of the $46$ Type C cases by automorphism order
(M\"oller-Scheithauer 2023, Table 3):
\begin{itemize}
\item $n = 2$: $15$ cases, e.g.\ $A_1^{24}$ deep hole orbifold;
\item $n = 3$: $7$ cases, including the $A_2^{12}$ / Coxeter-Todd
above;
\item $n = 4$: $7$ cases, e.g.\ $A_3^8$-type;
\item $n = 5$: $3$ cases;
\item $n = 6$: $6$ cases (composite orders include mixed-parity
contributions);
\item $n = 7$: $1$ case;
\item $n = 8$: $2$ cases;
\item $n = 10, 11, 12, 13, 14, 15, \ldots, 46$: $5$ cases total.
\end{itemize}
Each entry has a canonical orbifold presentation from
van~Ekeren--M\"oller--Scheithauer 2020 and a level-matching
verification in the cyclic-orbifold tables (twisted vacuum conformal weights, all lying in
$\frac{1}{n}\mathbb{Z}$ by construction).
\end{remark}

%%%%%%%%%%%%%%%%%%%%%%%%%%%%%%%%%%%%%%%%%%%%%%%%%%%%%%%%%%%%%%%%%%%%%%%%%%%%%
\section{Proof of the descent criterion}
\label{sec:schellekens-71-proof-assembly}
%%%%%%%%%%%%%%%%%%%%%%%%%%%%%%%%%%%%%%%%%%%%%%%%%%%%%%%%%%%%%%%%%%%%%%%%%%%%%

\begin{proof}[Proof of Theorem~\ref{thm:schellekens-71-all-alpha-zero}]
(i) The Schellekens row surface falls into exactly one of the three
strata by Definition~\ref{def:schellekens-orbifold-type} and the
van~Ekeren--M\"oller--Scheithauer construction theorem. Total coverage is
$24_{\mathrm{Niemeier}} + 1_{V_1=0} + 46_{\mathrm{nonlat}} = 71$
as enumerated in
Remark~\ref{rem:schellekens-strata-counts}.

(ii) The anomaly datum is stratum-dependent. Type~A has no finite
orbifold to gauge, so Proposition~\ref{prop:schellekens-typeA-alpha-zero}
gives the vacuous class. Type~B is the Leech $\Z/2$ local datum of
Proposition~\ref{prop:schellekens-typeB-alpha-zero}. Type~C uses the
cyclic-orbifold VOA input of Proposition~\ref{prop:ems-level-matching}
plus the explicit local trivialization required in
Corollary~\ref{cor:schellekens-typeC-alpha-zero}. These mechanisms are
pairwise distinct.

(iii) Chain-level $E_3$-topological descent: in Type~A, $V_\Lambda$ is
the boundary of abelian holomorphic Chern-Simons with charge lattice
$\Lambda$ (Theorem~\ref*{thm:E3-topological-km}, abelian case),
so chain-level $E_3$-topological structure is immediate. In Type~B,
Chapter~\ref*{ch:monster-chain-level-e3-top-platonic}
Theorem~\ref*{thm:monster-chain-level-e3-top} supplies the conditional
finite-orbifold descent for $V^\natural$. In Type~C, the charge
lattice is $\Lambda_{\mathrm{Leech}}$ and abelian holomorphic CS is
the parent HT model; the finite cyclic quotient inherits the
$E_3$ operations only after the BV anomaly has been trivialized and
finite-orbifold BV descent applies. Lattice automorphisms preserve the
abelian stress tensor $T = \frac{1}{2} :\sum J^a J^a:$, so the
topologisation is compatible with the action once the descent datum is
present.
\end{proof}

%%%%%%%%%%%%%%%%%%%%%%%%%%%%%%%%%%%%%%%%%%%%%%%%%%%%%%%%%%%%%%%%%%%%%%%%%%%%%
\section{Conditional chain-level $E_3$-topological descent in the Universal
Holography functor image}
\label{sec:schellekens-e3-uhf-image}
%%%%%%%%%%%%%%%%%%%%%%%%%%%%%%%%%%%%%%%%%%%%%%%%%%%%%%%%%%%%%%%%%%%%%%%%%%%%%

\begin{corollary}[Universal Holography functor on Schellekens $71$;
\ClaimStatusConditional]
\label{cor:schellekens-71-uhf-image}
The Universal Holography functor
$\Phi_{\mathrm{hol}}\colon \mathrm{ChirAlg}^{\omega,\mathrm{BL}}_X
\to \mathrm{HT\text{-}QFT}_{X \times \mathbb{R}}$
is defined on a Schellekens representative $W$ only after the
stratum-specific hypotheses of
Theorem~\ref{thm:schellekens-71-all-alpha-zero} are supplied. Under
those hypotheses, the descended HT factorisation algebra with boundary
$W$ is a chain-level HT-QFT on $X \times \mathbb{R}$ for every smooth
projective curve $X$.
\end{corollary}

\begin{proof}
Theorem~\ref{thm:schellekens-71-all-alpha-zero} provides chain-level
$E_3$-topological structure after the VOA presentation, anomaly
trivialization, and BV descent hypotheses are met. The Universal
Holography functor is defined on
chiral algebras with conformal vector at non-critical level
admitting chain-level $E_3$-topological refinement
(Theorem~\ref*{thm:universal-holography-functor}, hypotheses).
Schellekens VOAs are at $c = 24$ (non-critical for the present
topologisation theorem), but the conformal-vector source is
stratified: Type~A uses the lattice VOA conformal vector; Type~B uses
the FLM Moonshine conformal vector inherited through the Leech
$\mathbb Z/2$ orbifold; Type~C uses the conformal vector preserved by
the cyclic orbifold construction. Thus the functor applies
entry-by-entry only on the verified finite-orbifold descent locus. This
does not alter the generic class-$M$ original-complex restriction.
\end{proof}

\begin{remark}[Schellekens scope in the holography image]
\label{rem:uhf-class-coverage-update}
The Universal Holography functor has original-complex inputs on
classes G, L, C (via Theorem~\ref*{thm:E3-topological-km},
Theorem~\ref*{thm:E3-topological-DS-general}, and
Theorem~\ref*{thm:E3-topological-free-PVA}) and on $V^\natural$ via
the special Leech $\Z/2$ orbifold only under the local anomaly and BV
descent hypotheses of
Chapter~\ref*{ch:monster-chain-level-e3-top-platonic}. For the
Schellekens surface, Type~A is abelian, Type~B is the Monster
finite-orbifold criterion, and Type~C is cyclic-orbifold input plus a
local BV anomaly trivialization. This is a statement about the
holomorphic $c=24$ orbifold landscape; it does not change the generic
class-$M$ weight-completed ambient recorded in
Corollary~\ref*{cor:universal-holography-class-M}.
\end{remark}

%%%%%%%%%%%%%%%%%%%%%%%%%%%%%%%%%%%%%%%%%%%%%%%%%%%%%%%%%%%%%%%%%%%%%%%%%%%%%
\section{Relation to the Monster orbifold case}
\label{sec:schellekens-71-vs-monster}
%%%%%%%%%%%%%%%%%%%%%%%%%%%%%%%%%%%%%%%%%%%%%%%%%%%%%%%%%%%%%%%%%%%%%%%%%%%%%

The Monster $V^\natural$ is a special case (Type~B) of the present
criterion, distinguished by the shortcut
$\Lambda^\sigma = 0$. The Monster route is a finite-orbifold
criterion with a required local anomaly datum; Types~A and~C use
different inputs.

\begin{remark}[Three genuinely distinct mechanisms]
\label{rem:schellekens-three-mechanisms}
Type~A: no orbifold, anomaly class vacuous. The ``mechanism'' is that
there is nothing to compute; the underlying lattice is even
unimodular and the VOA is already a boundary of abelian holomorphic
CS.

Type~B: $\Lambda^\sigma = 0$ is the lattice input for the Leech
$\Z/2$ calculation. The HT descent still uses the local sign
$\alpha_{\mathrm{orb}}(\sigma,\sigma,\sigma)=+1$ and finite-orbifold
BV descent.

Type~C: $\Lambda^\sigma \ne 0$, so the $H^3$-cocycle restriction is a
genuine cohomology class. The cyclic-orbifold level-matching criterion
is the VOA condition; the HT claim requires the corresponding local
BV trivialization.

The three mechanisms are disjoint. No universal
``$\Lambda^\sigma$-triviality + even unimodularity'' argument, and no
Schellekens-list classification theorem alone, covers the HT descent
for the entire list.
\end{remark}

%%%%%%%%%%%%%%%%%%%%%%%%%%%%%%%%%%%%%%%%%%%%%%%%%%%%%%%%%%%%%%%%%%%%%%%%%%%%%
\section{Open directions and refinements}
\label{sec:schellekens-71-open-directions}
%%%%%%%%%%%%%%%%%%%%%%%%%%%%%%%%%%%%%%%%%%%%%%%%%%%%%%%%%%%%%%%%%%%%%%%%%%%%%

\begin{remark}[Beyond $c = 24$: arithmetic constraint]
\label{rem:schellekens-beyond-c24}
The Schellekens classification is specific to $c = 24$: the
requirement of strong rationality + holomorphy + central charge exactly
$24$ forces $V_1$ into the classified finite list. Extending the
present analysis to higher central charges $c \in 8\mathbb{Z}$ (the
values at which even unimodular lattices exist in positive definite
rank: $8, 16, 24, 32, 40, \ldots$) lies outside the present theorem:
at $c = 8$
there is one lattice VOA ($V_{E_8}$, with $V_1 = \mathfrak{e}_8$ at
level~$1$); at $c = 16$ there are two ($V_{E_8 \oplus E_8}$ and
$V_{D_{16}^+}$); at $c = 24$ the classified $71$; at $c = 32$ the
count is conjectural (millions of even unimodular lattices of
rank~$32$).
\end{remark}

\begin{remark}[Framed-VOA cross-check]
\label{rem:schellekens-framed-voa-crosscheck}
An alternative construction of the Schellekens $71$ via
framed-VOAs (Miyamoto 2004; Lam-Shimakura 2015) gives a different
presentation of each entry as an extension of the Virasoro tensor
product $\bigotimes L(c = 1/2, 0)^{\otimes 48}$. The framed-VOA
presentation gives an alternative VOA route toward anomaly analysis
via the $\mathbb{Z}/2$-simple-current extension theory of
Dong-Li-Mason 2000. It is a cross-check for the cyclic presentation
and simple-current phases; it does not replace the finite-orbifold BV
trivialization required in
Theorem~\ref{thm:schellekens-71-all-alpha-zero}.
\end{remark}

%%%%%%%%%%%%%%%%%%%%%%%%%%%%%%%%%%%%%%%%%%%%%%%%%%%%%%%%%%%%%%%%%%%%%%%%%%%%%
\section{Summary}
\label{sec:schellekens-71-summary}
%%%%%%%%%%%%%%%%%%%%%%%%%%%%%%%%%%%%%%%%%%%%%%%%%%%%%%%%%%%%%%%%%%%%%%%%%%%%%

\begin{itemize}
\item The $71$ Schellekens holomorphic $c = 24$ rows enter the HT
descent theorem through the conditional criterion
Theorem~\ref{thm:schellekens-71-all-alpha-zero}.
\item The inputs stratify as Type~A ($24$: trivial orbifold),
Type~B ($1$: Leech $\mathbb{Z}/2$, $\Lambda^\sigma = 0$ plus
local $\alpha_{\mathrm{orb}}=0$ datum), and Type~C ($46$: Leech
$\mathbb{Z}/n$ with $\Lambda^\sigma \ne 0$, cyclic-orbifold level
matching plus local BV trivialization).
\item The $\mathbb{Z}/3$ Coxeter--Todd calculation appears in
Proposition~\ref{prop:z3-coxeter-todd-alpha-zero-explicit}; every
other Type~C case has the same formal shape with different
$\Lambda^\sigma$ and different twisted-vacuum conformal weight.
\item Chain-level $E_3$-topological structure descends stratum by
stratum only after the abelian HT model, anomaly trivialization, and
finite-orbifold BV descent are supplied
(Corollary~\ref{cor:schellekens-71-uhf-image}).
\item The three mechanisms are genuinely distinct; the
$\Lambda^\sigma = 0$ shortcut of the Monster covers Type~B, not
Types~A and~C.
\item The stratified criterion preserves
$71 = 24_{\mathrm{Niemeier}} + 1_{V_1=0} + 46_{\mathrm{nonlat}}$.
\end{itemize}
